%% =====================================================================
%%  B3-T-04-03 -- what B3 contributes to "Vanity as a Handle"
%%  -------------------------------------------------------------------
%%  LAYER A: APPEND-ONLY. Written once, then left alone. Material found
%%  only here sits in the `unique` slot and is protected by rule R10.
%% =====================================================================
%% @kind: source
%% @id: B3-T-04-03
%% @book: B3
%% @topic: T-04-03
%% @locator: Law 1, pp. 1--7; Law 39
%% @status: distilled
%% @unique: yes
%% @integrated: yes
%% @updated: 2026-09-19

\dossier{B3}{T-04-03}{\bookshort{B3}}{Law 1, pp. 1--7; Law 39}

\begin{distilled}
  \item Displaying more talent than the person above you is described as inspiring fear and insecurity rather than admiration; the instruction is to make superiors appear more brilliant than they are.
  \item The reaction to a perceived slight is out of all proportion and there is no sanity behind it, so the counterparty should not be reasoned with.
  \item Finding the thread in an opponent's vanity and rattling it is presented as the way to take control of an exchange while staying calm yourself.
\end{distilled}

\begin{unique}
  \item The specific mechanism by which visible competence in a junior threatens a senior --- it is read as an intention to displace --- and the resulting instruction to make the superior's brilliance visible rather than your own.
\end{unique}
